% Generated by scripts/assemble.py; corrections belong in source/corrections.json.
\examhead{23}{中国矿业大学 2025 年数学分析试卷}{27}
% p028-b01; source PDF page 28
\par\Needspace{5\baselineskip}
\blockanchor{p028-b01}
\majorquestion{一}(15 分) 已知函数 $f(x)$ 是区间 $[a, b]$ 上的单调递增函数, 且 $f(a) > a$, $f(b) < b$. 证明: 存在一个数 $\xi \in (a, b)$, 使得 $f(\xi) = \xi$.\par
% p028-b02; source PDF page 28
\par\Needspace{5\baselineskip}
\blockanchor{p028-b02}
\majorquestion{二}（15分）已知函数 $f(x) = \frac{2x}{1 + x^2}$，求 $f(x)$ 的单调区间，极值点，凹凸区间，渐近线，并作出大致草图.\par
% p028-b03; source PDF page 28
\par\Needspace{5\baselineskip}
\blockanchor{p028-b03}
\majorquestion{三}(15 分) 已知 $f(x)$ 是 $[a, b]$ 上的连续函数, $x_1 < x_2 < \dots < x_n$ 是 $[a, b]$ 上的点.\par
% p028-b04; source PDF page 28
\blockanchor{p028-b04}
(1) (5 分) 叙述闭区间上的连续函数介值定理.\par
% p028-b05; source PDF page 28
\blockanchor{p028-b05}
(2) (10 分) 证明: 存在 $\xi \in [x_1, x_n]$, 使得 $f(\xi) = \sum_{j=1}^{n} \frac{f(x_j)}{n}$.\par
% p028-b06; source PDF page 28
\par\Needspace{5\baselineskip}
\blockanchor{p028-b06}
\majorquestion{四}(15 分) 求下列不定积分:\par
% p028-b07; source PDF page 28
\blockanchor{p028-b07}
(1) (5 分)  $\int \frac{1}{e^{x} - e^{-x}} dx.$\par
% p028-b08; source PDF page 28
\blockanchor{p028-b08}
(2) (10 分)  $\int \frac{1}{x + \sqrt{x^{2} - x + 1}} dx.$\par
% p028-b09; source PDF page 28
\par\Needspace{5\baselineskip}
\blockanchor{p028-b09}
\majorquestion{五}(15 分) 已知 $f(x)$ 在 $[a, b]$ 上可积.\par
% p028-b10; source PDF page 28
\par\Needspace{6\baselineskip}
\blockanchor{p028-b10}
(1) (5 分) 若 $f(x)$ 是 $[a, b]$ 上的连续非负函数, 且存在 $x_0 \in [a, b]$, 使得 $\lim_{x \to x_0} f(x) = A > 0$. 证明:\par
\nopagebreak[4]
% p028-b11; source PDF page 28
\blockanchor{p028-b11}
\displayformula{\int_ {a} ^ {b} f (x) \mathrm{d} x > 0.}
% p028-b12; source PDF page 28
\blockanchor{p028-b12}
(2) (10 分) 若  $\int_{a}^{b} f(x) \, dx > 0$ , 证明: 存在闭区间  $[\alpha, \beta] \subset [a, b]$ , 使得  $f(x)$  在  $[\alpha, \beta]$  上满足  $f(x) > 0$ .
\majorquestion{六}(15 分) 已知  $a_{n} = \int_{0}^{\frac{\pi}{4}} (\tan x)^{n} \, dx$ ，其中  $n = 1, 2, \cdots$ .\par
% p028-b13; source PDF page 28
\blockanchor{p028-b13}
(1) (5 分) 求级数 $\sum_{n=1}^{\infty} \frac{a_n + a_{n+2}}{n}$ 的和.\par
% p028-b14; source PDF page 28
\blockanchor{p028-b14}
(2) (10 分) 证明: 对任意常数 $p > 0$, 数项级数 $\sum_{n=1}^{\infty} \frac{a_n}{n^p}$ 都收敛.\par
% p028-b15; source PDF page 28
\par\Needspace{5\baselineskip}
\blockanchor{p028-b15}
\majorquestion{七}(15 分) 求重积分 $I = \iint_{D} \frac{3x}{y^2 + xy^3} \, \mathrm{d}x \, \mathrm{d}y$, 其中 $D = \{(x, y) \in \mathbb{R}^2 \mid 1 \leq xy \leq 3, x \leq y^2 \leq 3x\}$.\par
% p028-b16; source PDF page 28
\par\Needspace{5\baselineskip}
\blockanchor{p028-b16}
\majorquestion{八}(15 分) 已知 $f(x, y)$ 在 $P_0(x_0, y_0)$ 的某邻域 $U(P_0)$ 上有二阶连续偏导数，且 $f_x(x_0, y_0) = 0$，$f_{xx}(x_0, y_0) > 0$. 证明：存在关于 $y_0$ 的某邻域 $V(y_0)$，对任意 $y \in V(y_0)$，$f(x, y)$ 作为关于 $x$ 的函数存在极小值 $g(y)$，且 $g'(y_0) = f_y(x_0, y_0)$.\par
% p028-b17; source PDF page 28
\par\Needspace{5\baselineskip}
\blockanchor{p028-b17}
\majorquestion{九}(15 分) 已知 $u(x, y)$ 在 $\mathbb{R}^2$ 上有二阶连续偏导数, 若 $u_{xx} + u_{yy} = 0$ 对任意 $(x, y) \in \mathbb{R}^2$ 都成立, 则称 $u(x, y)$ 为 $\mathbb{R}^2$ 上的调和函数. 证明: $u(x, y)$ 为 $\mathbb{R}^2$ 上的调和函数的充要条件是 $\oint_C \frac{\partial u}{\partial n} \mathrm{d}s = 0$, 其中 $C$ 为 $\mathbb{R}^2$ 上任意光滑闭曲线, $\pmb{n}$ 为 $u(x, y)$ 沿 $C$ 的单位外法线方向.\par
% p028-b18; source PDF page 28
\par\Needspace{5\baselineskip}
\blockanchor{p028-b18}
\majorquestion{十}(15 分) 已知欧拉积分 $\int_0^{+\infty} e^{-x^2} \, \mathrm{d}x = \frac{\sqrt{\pi}}{2}$, 求含参量积分 $I(a) = \int_0^{+\infty} e^{-\left(x^2 + \frac{a^2}{x^2}\right)} \, \mathrm{d}x$ 的值, 其中 $a > 0$.\par
